% !TEX root = ../ComputationalOTFnT.tex

\chapter{最適輸送のエントロピー正則化}
\label{c-entropic}

この章では、最適輸送のカントロビッチ定式化とその多くの一般化に対する近似解を求めるための数値スキームの族を導入する。これはもとの問題にエントロピー正則化ペナルティを加えることで動作する。この正則化にはいくつかの重要な利点があり、それらを総合すると非常に有用なツールとなる：正則化された問題の最小化は単純な交互最小化スキームで解くことができる；そのスキームは単純な行列-ベクトル積の反復に帰着し、GPU上での実行に特に適している；一部の応用では、これらの行列-ベクトル積は$n\times m$のコスト行列を格納する必要がなく、カーネル評価へのアクセスのみを必要とする；大規模な測度群が同じ台を共有する場合、これらの行列-ベクトル積はすべて行列-行列積として実行でき、大幅な高速化が得られる；得られる近似距離は入力ヒストグラムの重みとDirac測度の位置に関して滑らかであり、自動微分を用いて微分可能である。

%%%%%%%%%%%%%%%%%%%%%%%%%%%%%%%%%%%%%%%%%%%
\section{エントロピー正則化}

カップリング行列の離散エントロピーは次のように定義される
\eql{\label{eq-discr-entropy}
	\HD(\P) \eqdef -\sum_{i,j} \P_{i,j} (\log(\P_{i,j})-1),
}
ベクトルに対しても同様の定義が成り立ち、$\a_j$のいずれかの成分が0または負の場合は$\HD(\a) = -\infty$とする規約を設ける。関数$\HD$は1-強凹関数である。なぜならそのHesse行列は$\partial^2 \HD(P)=-\diag(1/\P_{i,j})$であり$\P_{i,j} \leq 1$だからである。最適輸送のエントロピー正則化のアイデアは、$-\HD$を正則化関数として用い、もとの輸送問題~\eqref{eq-mk-discr}の近似解を得ることである：
\eql{\label{eq-regularized-discr}
	\MKD_\C^\varepsilon(\a,\b) \eqdef
	\umin{\P \in \CouplingsD(\a,\b)}
		\dotp{\P}{\C} - \varepsilon \HD(\P).
}
目的関数は$\varepsilon$-強凸関数であるため、問題(\ref{eq-regularized-discr})は一意の最適解を持つ。最適輸送問題をエントロピー項で正則化するというアイデアは、交通理論におけるモデリングのアイデアにまで遡ることができる~\citep{wilson1969use}：ネットワークにおける実際の交通パターンは、最適輸送問題の解によって予測されるものとは一致しない。実際、後者は~\eqref{eq-mk-discr}の最適カップリングのスパース性の結果として少数のルートに依存する傾向があるのに対し、前者はより拡散的である。このスパース性を緩和するために、交通の研究者は「重力」モデル~\citep{erlander1980optimal}と呼ばれるモデルを提案し、周辺分布と輸送コストが与えられた場合に、より「ぼやけた」交通予測を形成できるようにした。

図~\ref{fig-impact-eps}は、単体$\simplex_3$（したがって2次元の三角形として可視化できる）上の線形計画法を正則化するエントロピーの効果を示している。エントロピーがもとのLP解を三角形の境界から押し離す様子に注目されたい。最適な$\P_\varepsilon$は徐々に三角形の「エントロピー的中心」に向かって移動する。これは以下の命題でさらに詳しく述べる。正則化された問題の解がもとの線形計画法の最適解に収束することは~\citet{CominettiAsympt}によって研究されており、精密な漸近展開が与えられている。


\begin{figure}[h!]
\centering
\includegraphics[width=.7\linewidth]{entropic/simplex}
\caption{\label{fig-impact-eps}
単体上の線形関数の最適化に対する$\varepsilon$の影響、様々な$\varepsilon$に対して$\P_\varepsilon = \argmin_{\P \in \simplex_3} \dotp{\C}{\P}-\varepsilon\HD(\P)$を解く。
}
\end{figure}


\begin{prop}[$\varepsilon$に関する収束]\label{prop-convergence-eps}
~\eqref{eq-regularized-discr}の一意の解$\P_\varepsilon$は、カントロビッチ問題のすべての最適解の集合の中で最大エントロピーを持つ最適解に収束する。すなわち
\eql{\label{eq-entropy-conv-1}
	\P_\varepsilon \overset{\varepsilon \rightarrow 0}{\longrightarrow}
	\uargmin{\P} \enscond{ -\HD(\P) }{
		\P \in \CouplingsD(\a,\b), \dotp{\P}{\C} = \MKD_\C(\a,\b),
	}
}
したがって特に
\eq{
	\MKD_\C^\varepsilon(\a,\b) \overset{\varepsilon \rightarrow 0}{\longrightarrow} \MKD_\C(\a,\b).
}
また
\eql{\label{eq-entropy-conv-2}
	\P_\varepsilon \overset{\varepsilon \rightarrow \infty}{\longrightarrow}
	\a \otimes \b = \a \transp{\b} = (\a_i \b_j)_{i,j}.
}
\end{prop}
\begin{proof}
	 $\varepsilon_\ell \rightarrow 0$かつ$\varepsilon_\ell > 0$である列$(\varepsilon_\ell)_\ell$を考える。
 	$\varepsilon=\varepsilon_\ell$のときの~\eqref{eq-regularized-discr}の解を$\P_\ell$と表す。
	%
	$\CouplingsD(\a,\b)$は有界であるため、（簡単のため添字を付け替えない）部分列を抽出して$\P_\ell \rightarrow \P^\star$とすることができる。$\CouplingsD(\a,\b)$は閉集合であるから、$\P^\star \in \CouplingsD(\a,\b)$である。$\dotp{\C}{\P} = \MKD_\C(\a,\b)$を満たす任意の$\P$を考える。$\P$と$\P_\ell$がそれぞれの最適化問題（$\varepsilon=0$と$\varepsilon=\varepsilon_\ell$）に対して最適であることから、
 	\eql{\label{eq-proof-gamma-conv}
 		0 \leq \dotp{\C}{\P_\ell} - \dotp{\C}{\P} \leq \varepsilon_\ell ( \HD(\P_\ell)-\HD(\P) ).
 	}
 	$\HD$は連続であるから、この式で$\ell \rightarrow +\infty$の極限をとると
 	$\dotp{\C}{\P^\star} = \dotp{\C}{\P}$が得られ、$\P^\star$は~\eqref{eq-entropy-conv-1}の実行可能点である。さらに、~\eqref{eq-proof-gamma-conv}で$\varepsilon_\ell$で割って極限をとると
 	$\HD(\P) \leq \HD(\P^\star)$が得られ、$\P^\star$が~\eqref{eq-entropy-conv-1}の解であることが示される。$-\HD$の狭義凸性によりこの計画の解$\P_0^\star$は一意であるから、$\P^\star = \P_0^\star$であり、列全体が収束する。
	%
	$\varepsilon\rightarrow+\infty$の極限においても同様の証明により、むしろ問題
	\eq{
		\umin{\P \in \CouplingsD(\a,\b)} -\HD(\P),
	}
	を考えるべきであり、その解は$\a \otimes \b$である。
\end{proof}

式~\eqref{eq-entropy-conv-1}は、小さな正則化$\varepsilon$に対して、解が最大エントロピー最適輸送カップリングに収束することを述べている。
%
対照的に、~\eqref{eq-entropy-conv-2}は、大きな正則化に対して、解が規定された周辺分布$\a,\b$の間の最大エントロピーカップリング、すなわち$\a,\b$に従って分布する2つの独立な確率変数の同時確率に収束することを示している。
%
この収束のより精密な分析は~\citet{CominettiAsympt}によって行われており、$\varepsilon=0$（resp., $\varepsilon=+\infty$）の近傍における$\varepsilon$（resp., $1/\varepsilon$）の1次展開を含む。
%
図~\ref{fig-entropic-densities}と図~\ref{fig-entropic}はこれら2つの収束の効果を視覚的に示している。重要な知見は、$\varepsilon$が増加するにつれて、最適カップリングはますますスパースでなくなり（所定の閾値より大きい成分を持つという意味で）、その結果、計算アルゴリズムの高速化（\S\ref{sec-sinkhorn}で検討）と統計的収束の高速化（\S\ref{sec-entropy-ot-mmd}で示される）の両方の効果をもたらすということである。

\newcommand{\myfigSinkEps}[1]{%
\imgBox{\includegraphics[width=.23\linewidth]{sinkhorn-eps/evol-img-#1}}}
\newcommand{\myfigSinkEpsV}[1]{\includegraphics[width=.24\linewidth]{sinkhorn-eps/evol-3d-#1}}

\begin{figure}[h!]
\centering
\begin{tabular}{@{}c@{}c@{}c@{}c@{}}
\myfigSinkEps{1} &
\myfigSinkEps{3} &
\myfigSinkEps{4} &
\myfigSinkEps{5} \\
\myfigSinkEpsV{1} &
\myfigSinkEpsV{3} &
\myfigSinkEpsV{4} &
\myfigSinkEpsV{5} \\
$\varepsilon=10$ &
$\varepsilon=1$ &
$\varepsilon=10^{-1}$ &
$\varepsilon=10^{-2}$
\end{tabular}
\caption{\label{fig-entropic-densities}
カップリングに対する$\varepsilon$の影響、命題~\ref{prop-convergence-eps}の図解。
%
上段：2つの1次元密度間。下段：同じ点数$n=m$の2つの2次元離散経験密度間（最適な$(\P_{i,j})_{i,j}$の小さな閾値以上の成分のみが$x_i$と$y_j$の間の線分として表示されている）。
}
\end{figure}


\begin{figure}[h!]
\centering
\includegraphics[width=\linewidth]{entropic/matching-2d}
\caption{\label{fig-entropic}
同じ点数$n=m$の2つの2次元離散経験密度間のカップリングに対する$\varepsilon$の影響（最適な$(\P_{i,j})_{i,j}$の小さな閾値以上の成分のみが$x_i$と$y_j$の間の線分として表示されている）。
}
\end{figure}

カップリング間のKullback-Leibler・ダイバージェンスを次のように定義すると
\eql{\label{eq-kl-defn}
	\KLD(\P|\K) \eqdef \sum_{i,j}  \P_{i,j} \log\pa{\frac{\P_{i,j}}{\K_{i,j}}} - \P_{i,j} + \K_{i,j},
}
~\eqref{eq-regularized-discr}の一意の解$\P_\varepsilon$は、コスト行列$\C$に関連するGibbsカーネルの$\CouplingsD(\a,\b)$への射影として得られる
\eq{
	\K_{i,j} \eqdef e^{-\frac{\C_{i,j}}{\varepsilon}}.
}
実際、上の定義を用いると
\eql{\label{eq-kl-proj}
	\P_\varepsilon = \Proj_{\CouplingsD(\a,\b)}^\KLD(\K) \eqdef \uargmin{\P \in \CouplingsD(\a,\b)} \KLD(\P|\K).
}




%%%%%
\begin{rem1}{離散測度間のエントロピー正則化}\label{rem:entrop-reg}
~\eqref{eq-discr-meas}の形の離散測度に対して、正則化された輸送の定義は自然に次のように拡張される
\eql{\label{eq-entropic-disc-meas}
	\MK_\c^\varepsilon(\al,\be) \eqdef \MKD_\C^\varepsilon(\a,\b),
}
コスト$\C_{i,j}=\c(x_i,y_j)$であり、入力測度を支持する位置$(x_i,y_j)$への依存性を強調している。
\end{rem1}
%%%%%


%%%%%
\begin{rem2}{一般的な定式化}
離散エントロピーを直積測度$\d\al\otimes\d\be(x,y) \eqdef \d\al(x)\d\be(y)$に関する相対エントロピーで置き換えることで任意の測度を考えることができ、~\eqref{eq-mk-generic}に対する正則化された対応物として次を提案する
\eql{\label{eq-entropic-generic}
	\MK_\c^\varepsilon(\al,\be) \eqdef
	\umin{\pi \in \Couplings(\al,\be)}
		\int_{\X \times \Y} c(x,y) \d\pi(x,y) + \varepsilon \KL(\pi|\al\otimes\be),
}
ここで相対エントロピーは離散Kullback-Leibler・ダイバージェンス~\eqref{eq-kl-defn}の一般化である
\eql{\label{eq-defn-rel-entropy}
\begin{aligned}	 \KL(\pi|\xi) \eqdef \int_{\X \times \Y} \log\Big( \frac{\d \pi}{\d\xi}(x,y) \Big) \d\pi(x,y)\\
	  + \int_{\X \times \Y} (\d\xi(x,y)-\d\pi(x,y)),
	  \end{aligned}
}
$\pi$が$\xi$に関する密度$\frac{\d \pi}{\d\xi}$を持たない場合は$\KL(\pi|\xi)=+\infty$と約束する。
%
~\eqref{eq-entropic-generic}でエントロピー正則化項$\KL(\cdot|\al\otimes\be)$を定義するために選んだ基準測度$\al\otimes\be$は特別な役割を果たさないことを認識することが重要である；以下の命題が述べるように、その台のみが重要である。

\begin{prop}
	任意の$\pi \in \Couplings(\al,\be)$に対して、$(\al,\be)$と同じゼロ測度集合を持つ任意の$(\al',\be')$（したがって互いに密度を持つ）に対して
	\eq{
		\KL(\pi|\al\otimes\be) = \KL(\pi|\al'\otimes\be') - \KL(\al\otimes\be|\al'\otimes\be').
	}
\end{prop}
\todoK{
\begin{proof}
	\todo{Prove it}
\end{proof}
}

この命題は、~\eqref{eq-entropic-generic}において$\KL(\cdot|\al\otimes\be)$の代わりに$\KL(\cdot|\al'\otimes\be')$を選んでも同じ解が得られることを示している。


式~\eqref{eq-entropic-generic}は射影問題として再定式化できる
\eql{\label{eq-entropic-generic-proj}
	\umin{\pi \in \Couplings(\al,\be)} \KL(\pi|\Kk),
}
ここで$\Kk$はGibbs分布$\d\Kk(x,y) \eqdef e^{-\frac{c(x,y)}{\varepsilon}} \d\mu(x)\d\nu(y)$である。
%
この問題は「静的Schrödinger問題」~\citep{LeonardSchroedinger,RuschendorfThomsen}と呼ばれることが多い。これはもともとSchrödingerが統計物理学において考えた問題である~\citep{Schroedinger31}。
%
$\varepsilon \rightarrow 0$のとき、~\eqref{eq-entropic-generic-proj}の一意の解は~\eqref{eq-mk-generic}の最大エントロピー解に収束する；~\citep{leonard2012schrodinger,2017-carlier-SIMA}を参照。
%
\S\ref{sec-entropic-dynamic}では、2つの測度の点を結ぶ経路の空間上でのSchrödinger問題の代替的な「動的」定式化を詳しく述べる。
\end{rem2}
%%%%%

%%%%%
\begin{rem2}{相互エントロピー}
~\eqref{eq-ot-proba-interpretation}と同様に、確率変数を用いて\eqref{eq-entropic-generic}を次のように言い換えることができる
\eq{
	\MK_\c^\varepsilon(\al,\be) = \umin{(X,Y)} \enscond{ \EE_{(X,Y)}(c(X,Y)) + \epsilon \text{I}(X,Y) }{ X \sim \al, Y \sim \be },
}
ここで$\pi$を$(X,Y)$の分布として、$\text{I}(X,Y) \eqdef \KL(\pi|\al\otimes\be)$は2つの確率変数間のいわゆる相互情報量である。$\text{I}(X,Y) \geq 0$であり、$\text{I}(X,Y)=0$となるのは2つの確率変数が独立である場合に限る。
\end{rem2}

%%%%%
\begin{rem2}{独立性とカップリング}
	カップリング$\pi \in \Couplings(\al,\be)$は、$(\X,\Y)$上で定義された確率変数の組$(X,Y)$の分布を記述する。ここで$X$（resp., $Y$）は法則$\al$（resp., $\be$）に従う。
	%
	命題~\ref{prop-convergence-eps}は一般的な（必ずしも離散でない）測度に対しても成り立つ。したがって~\eqref{eq-entropic-generic}の解$\pi_\varepsilon$は$\varepsilon \rightarrow+\infty$でテンソル積カップリング$\al \otimes \be$に収束する。このカップリング$\al \otimes \be$は確率変数$(X,Y)$が独立であることに対応する。
	%
	対照的に、$\varepsilon \rightarrow 0$のとき、$\pi_\varepsilon$はOT問題~\eqref{eq-mk-generic}の解$\pi_0$に収束する。$\X=\Y=\RR^\dim$において、$\al$と$\be$がLebesgue測度に関する密度を持つ場合、注意~\ref{rem-exist-mongemap}で詳述するように、$\pi_0$は一意であり、全単射のMonge写像$\T : \RR^\dim \rightarrow \RR^\dim$のグラフ上に台を持つ。この場合、$Y=\T(X)$かつ$X=\T^{-1}(Y)$であるため、$(X,Y)$はある意味で完全に従属している。
	%
	$\dim=1$の単純な1次元の場合、$X$と$Y$の従属構造を視覚化する便利な方法は、同時分布$\pi$に付随するコピュラ$\xi_\pi$を用いることである。~\eqref{eq-cumul-defn}で定義された累積分布関数は、カップリングに対して次のように拡張される
	\eq{
		\foralls (x,y) \in \RR^2, \quad
		\cumul{\pi}(x,y) \eqdef \int_{-\infty}^x \int_{-\infty}^y \d\pi.
	}
	コピュラは次のように定義される
	\eq{
		\foralls (s,t) \in [0,1]^2, \quad
		\xi_{\pi}(s,t) \eqdef \cumul{\pi}( \cumul{\al}^{-1}(s),\cumul{\be}^{-1}(t) ),
	}
	ここで累積分布関数の擬似逆関数は~\eqref{eq-pseudo-inv-cum}で定義される。
	%
	独立な変数の場合、$\varepsilon=+\infty$、すなわち$\pi=\al\otimes \be$のとき、$\xi_{\pi_{+\infty}}(s,t)=st$となる。
	%
	対照的に、完全に従属する変数の場合、$\varepsilon=+\infty$のとき、$\xi_{\pi_0}(s,t)=\min(s,t)$となる。
	%
	図~\ref{fig-entropic-copula}は、エントロピー正則化がこれら2つの極端な場合を補間するコピュラ$\xi_{\pi_\varepsilon}$をどのように生成するかを示している。
\end{rem2}
%%%%%

% G B D H
\newcommand{\myfigSinkCopul}[1]{\imgBox{\includegraphics[width=.185\linewidth,trim=65 55 45 40,clip]{sinkhorn-copula/evol-#1}}}
\begin{figure}[h!]
\centering
\begin{tabular}{@{}c@{\hspace{5mm}}c@{}}
\includegraphics[width=.4\linewidth]{sinkhorn-copula/input-1} &
\includegraphics[width=.4\linewidth]{sinkhorn-copula/input-2}\\
{\color{red}$\al$} & {\color{blue}$\be$}
\end{tabular}
%
\begin{tabular}{@{}c@{}c@{}c@{}c@{}c@{}}
\myfigSinkCopul{levelsets-1} &
\myfigSinkCopul{levelsets-2} &
\myfigSinkCopul{levelsets-3} &
\myfigSinkCopul{levelsets-4} &
\myfigSinkCopul{levelsets-5} \\
\myfigSinkCopul{copula-1} &
\myfigSinkCopul{copula-2} &
\myfigSinkCopul{copula-3} &
\myfigSinkCopul{copula-4} &
\myfigSinkCopul{copula-5} \\
{\color[rgb]{1,0,1} $\varepsilon=10$} &
{\color[rgb]{.75,.25,.75} $\varepsilon=1$} &
{\color[rgb]{.5,.5,.5} $\varepsilon=0.5 \cdot 10^{-1}$} &
{\color[rgb]{.25,.75,.25} $\varepsilon=10^{-1}$} &
{\color[rgb]{0,1,0} $\varepsilon=10^{-3}$}
\end{tabular}
\caption{\label{fig-entropic-copula}
上：~\eqref{eq-entropic-generic}の解$\pi_\varepsilon$の$\varepsilon$による変化。
下：コピュラ関数$\xi_{\pi_\varepsilon}$の変化。
}
\end{figure}


%%%%%%%%%%%%%%%%%%%%%%%%%%%%%%%%%%%%%%%%%%%
\section{Sinkhornアルゴリズムとその収束}
\label{sec-sinkhorn}

以下の命題は、~\eqref{eq-regularized-discr}の解が$n+m$個の変数でパラメータ化できる特定の形を持つことを示す。このパラメータ化は本質的に双対的である。なぜなら$\CouplingsD(\a,\b)$におけるカップリング$\P$は$nm$個の変数を持つが$n+m$個の制約を持つからである。

\begin{prop}\label{prop-regularized-primal}
~\eqref{eq-regularized-discr}の解は一意であり、次の形を持つ
\eql{\label{eq-scaling-form}
	\foralls (i,j) \in \range{n} \times \range{m}, \quad \P_{i,j} = \uD_i \K_{i,j} \vD_j
}
ここで$(\uD,\vD) \in \RR_+^n \times \RR_+^m$は2つの（未知の）スケーリング変数である。
\end{prop}

\begin{proof}
各周辺制約に対する2つの双対変数$\fD\in\RR^n,\gD\in\RR^m$を導入すると、~\eqref{eq-regularized-discr}のLagrangianは次のようになる
\eq{\label{eq-sinkhorn-lagrangian}
	\Lag(\P,\fD,\gD)= \dotp{\P}{\C} - \varepsilon \HD(\P) - \dotp{\fD}{\P\ones_m-\a}-\dotp{\gD}{\transp{\P}\ones_n-\b}.
}
一階条件から
$$
	\frac{\partial\Lag(\P,\fD,\gD)}{\partial \P_{i,j}}= \C_{i,j} + \varepsilon \log(\P_{i,j}) - \fD_i -\gD_j=0,
$$
が得られ、正則化問題に対する最適カップリング$\P$は$\P_{i,j}=e^{\fD_i/\varepsilon}e^{-\C_{i,j}/\varepsilon}e^{\gD_j/\varepsilon}$と表される。これは非負ベクトル$\uD$と$\vD$を用いて上で与えた形に書き換えることができる。
\end{proof}

%%%
\paragraph{行列スケーリングとしての正則化OT}

式~\eqref{eq-scaling-form}で示された最適解の因子分解は、行列形式で$\P=\diag(\uD)\K\diag(\vD)$と便利に書き換えることができる。
%
したがって変数$(\uD,\vD)$は、$\CouplingsD(\a,\b)$に固有の質量保存制約に対応する次の非線形方程式を満たさなければならない：
\eql{\label{eq-dualsinkhorn-constraints}
	\diag(\uD)\K\diag(\vD)\ones_m=\a,
	\qandq
	\diag(\vD)\K^\top \diag(\uD)\ones_n=\b.
}
これら2つの方程式はさらに簡略化できる。$\diag(\vD)\ones_m$は単に$\vD$であり、$\diag(\uD)$と$\K \vD$の積は
\eql{\label{eq-dualsinkhorn-constraints2}
	\uD \odot (\K \vD) = \a
	\qandq
	\vD \odot (\transp{\K}\uD) = \b,
}
ここで$\odot$はベクトルの成分ごとの積に対応する。この問題は数値解析のコミュニティでは行列スケーリング問題として知られている（~\citep{nemirovski1999complexity}およびその参考文献を参照）。
%
これらの方程式を扱う直観的な方法は、まず$\uD$を式~\eqref{eq-dualsinkhorn-constraints2}の左辺を満たすように修正し、次に$\vD$をその右辺を満たすように修正して、反復的に解くことである。これら2つの更新がSinkhornアルゴリズムを定義する
\eql{\label{eq-sinkhorn}
	\itt{\uD} \eqdef \frac{\a}{\K \it{\vD}}
	\qandq
	\itt{\vD} \eqdef \frac{\b}{\transp{\K}\itt{\uD}},
}
初期値は任意の正ベクトル$\init{\vD} = \ones_m$である。上で用いた2つのベクトル間の除算演算子は成分ごとと理解されたい。異なる初期化は$\uD,\vD$に対して異なる解につながる可能性がある。なぜなら$\uD,\vD$は乗法定数の不定性を持つからである（$\uD,\vD$が\eqref{eq-dualsinkhorn-constraints}を満たすならば、任意の$\lambda>0$に対して$\lambda\uD,\vD/\lambda$も満たす）。
%
しかしながら、これらの反復は収束し（反復射影を用いた正当化については注意~\ref{rem-iterative-projection}を、厳密な縮小性の結果については注意~\ref{rem-global-conv-sinkh}を参照）、すべて同じ最適カップリング$\diag(\uD)\K\diag(\vD)$に帰着する。
%
図~\ref{fig-sinkhorn-convergence}の上段は、Sinkhorn反復によって計算されるカップリング$\diag(\it{\uD})\K\diag(\it{\vD})$の変化を示している。Gibbsカーネル$\K$から~\eqref{eq-regularized-discr}を解く最適カップリングへ、質量を対角線から徐々にずらしながら収束していく様子がわかる。

\begin{rem}[歴史的展望] 反復~\eqref{eq-sinkhorn}は最初に~\citep{yule1912methods,kruithof}に登場した。後にIPFP（iterative proportional fitting procedure）~\citet{DemingStephanIPFP}およびRAS~\citep{bacharach1965estimating}法~\citep{ReviewSinkhorn}として知られるようになった。その収束の証明は~\citet{Sinkhorn64}に帰せられ、アルゴリズムの名前の由来となっている。このアルゴリズムは後に~\citet{Ruschendorf95}によって無限次元に拡張された。この正則化は経済学の分野で最適輸送問題の近似解を得るために使用され、重力モデル~\citep{wilson1969use,erlander1980optimal,erlander1990gravity}の名で知られていた。
% with later contributions~\citep{SinkhornKnopp67,Sinkhorn67}
~\citet{kosowsky1994invisible}により、割当問題の場合、すなわち$\a=\b=\ones_n/n$の場合に「ソフト割当（softassign）」として再ブランド化され、より最近では~\citet{Galichon-Entropic}によって経済学におけるマッチング問題を解くために使用された。
%
この正則化は、~\citep{CuturiSinkhorn}に続いてデータサイエンス（機械学習、ビジョン、グラフィックス、画像処理を含む）において再び注目を集めた。彼らはSinkhornアルゴリズムが最適輸送に対する効率的でスケーラブルな近似を提供することを示した。これは複数のOT問題を同時に解く際のシームレスな並列化（特にGPU上；注意~\ref{rem-parallel}を参照）のおかげであり、またこの正則化された量は線形計画法の定式化とは異なり微分可能な損失関数を定義することも示した（\S\ref{sec-regularized-cost}を参照）。
%
Sinkhornアルゴリズムの拡張と一般化は無数に存在する（例えば\S\ref{sec-generalized}を参照）。例えば、$\a=\b$の場合、対称性を維持するために平均化された射影反復を用いることができる~\citep{knight2014symmetry}。
\end{rem}

\begin{figure}[h!]
\centering
\includegraphics[width=\linewidth]{entropic/sinkhorn-convergence}
\begin{tabular}{c}
\includegraphics[width=.5\linewidth]{entropic/sinkhorn-rates}\\
{\color{blue}$\varepsilon=10$} \quad {\color[rgb]{0.5,0,0.5}$\varepsilon=0.1$} \quad {\color{red}$\varepsilon=10^{-3}$}
\end{tabular}
\caption{\label{fig-sinkhorn-convergence}
上：$\X=[0,1]$上の1次元密度、$c(x,y)=|x-y|^2$、$\varepsilon=0.1$の場合の、Sinkhorn反復の$\ell$回目の反復で計算されるカップリング$\it{\pi_\varepsilon}=\diag(\it{\uD})\K\diag(\it{\vD})$の変化。
下：周辺制約違反$\log( \|\it{\pi_\varepsilon} \ones_m - \b \|_1 )$で測定したSinkhornの収束速度に対する$\varepsilon$の影響。
}
\end{figure}


%%%%%%%%%%
\begin{rem1}{全体の計算量}\label{rem-complexity-rounding}
	注意深い収束解析により（簡単のため$n=m$を仮定）、\citet{altschuler2017near}は$\varepsilon = \frac{4\log(n)}{\tau}$と設定した場合、$O(\norm{\C}_\infty^3 \log(n) \tau^{-3})$回のSinkhorn反復（有効なカップリング$\hat \P \in \CouplingsD(\a,\b)$を計算するための追加の丸めステップを伴う）で$\dotp{\hat \P}{\C} \leq \MKD_{\C}(\a,\b) + \tau$を保証するのに十分であることを示した。これはSinkhornが非正則化OT問題の$\tau$-近似解を$O(n^2\log(n) \tau^{-3})$演算で計算することを意味する。丸めスキームは、2つのベクトル$\uD\in\RR^n,\vD\in\RR^m$が与えられたとき、以下の更新を行うものである（\cite[Alg. 2]{altschuler2017near}）：
	$$\begin{aligned}
		\uD'& \eqdef \uD\odot\min\left(\frac{\a}{\uD\odot(\K\vD)},\ones_n\right),
		\vD' \eqdef \vD\odot \min\left(\frac{\b}{\vD\odot(\transp{\K}\uD')},\ones_n\right), \\
	\Delta_{\a}& \eqdef \a- \uD'\odot (\K \vD'), \Delta_{\b} \eqdef \b- \vD'\odot (\transp{\K} \uD), \\
	\hat\P& \eqdef \diag(\uD')\K\diag(\vD')+\Delta_{\a} \transp{(\Delta_{\b})}/\norm{\Delta_{\a}}_1.
	\end{aligned}$$
	これにより$\hat\P \in \CouplingsD(\a,\b)$である行列が得られ、$\hat\P$と$\diag(\uD)\K\diag(\vD)$の$1$-ノルムの差は$\diag(\uD)\K\diag(\vD)$の周辺違反によって制御される。すなわち
	$$ \norm{\hat\P-\diag(\uD)\K\diag(\vD)}_1 \leq \norm{\a-\uD\odot(\K\vD)}_1+\norm{\b-\vD\odot(\transp{\K}\uD)}_1.$$
この分野は活発であり、\citet{pmlr-v80-dvurechensky18a}による上記の結果の最近の改善がそれを示している。
\end{rem1}
%%%%%%%%%%

%%%%%%%
\begin{rem}[Sinkhorn反復の数値安定性]\label{rem-stability}
注意~\ref{rem-global-conv-sinkh}と\ref{rem-local-conv}で議論するように、Sinkhornアルゴリズムの収束は$\varepsilon\rightarrow 0$で悪化する。しかし数値計算上の実際では、より単純な理由でその減速が観察されることはまれである：カーネル$\K$の要素の一部がメモリ上に正の数として格納するには小さすぎてゼロになると、Sinkhornアルゴリズムはしばしば終了に失敗する。この結果、行列積$\K\vD$や$\transp{\K}\uD$の成分がますます小さくなってゼロになり、式~\eqref{eq-sinkhorn}のSinkhorn更新で0での除算が生じる。このような問題は、乗数$\uD$と$\vD$の計算を対数領域で行うことで部分的に解決できる。このアプローチは注意~\ref{rem-log-sinkh}で注意深く説明され、問題~\eqref{eq-regularized-discr}の双対の直接的な解法と関連している。
\end{rem}
%%%%%%%

%%%%%%%
\begin{rem}[反復射影との関係]\label{rem-iterative-projection}
\eq{
	\Cc^1_\a \eqdef \enscond{\P}{\P\ones_m=\a}
	\qandq
	\Cc^2_\b \eqdef \enscond{\P}{\transp{\P}\ones_m=\b}
}
と行と列の制約を表すと、$\CouplingsD(\a,\b) = \Cc^1_\a \cap \Cc^2_\b$となる。Bregman反復射影~\citep{bregman1967relaxation}を用いることができる
\eql{\label{eq-kl-sinkh-proj}
	\itt{\P} \eqdef \Proj_{\Cc^1_\a}^{\KLD}(\it{\P})
	\qandq
	\ittt{\P} \eqdef \Proj_{\Cc^2_\b}^{\KLD}(\itt{\P}).
}
集合$\Cc^1_\a$と$\Cc^2_\b$はアフィンであるため、これらの反復は~\eqref{eq-kl-proj}の解に収束することが知られている；~\citep{bregman1967relaxation}を参照。これらの反復はSinkhorn反復~\eqref{eq-sinkhorn}と等価であり、
\eq{\label{eq-sink-matrix}\P^{(2\ell)} \eqdef \diag(\it{\uD}) \K \diag(\it{\vD}),}
と定義すると
\begin{align*}
	\P^{(2\ell+1)} &\eqdef \diag(\itt{\uD}) \K \diag(\it{\vD}) \\
	\qandq
	\P^{(2\ell+2)} &\eqdef \diag(\itt{\uD}) \K \diag(\itt{\vD}).
\end{align*}
しかし実際には、スケーリングベクトルの操作とGibbsカーネルとの乗算のみを必要とする~\eqref{eq-sinkhorn}を用いることが好ましい。これはしばしば高速化できるからである（以下の注意~\ref{rem-separable}と~\ref{rem-geod-heat}を参照）。
\end{rem}
%%%%%%%

%%%%%%%
\begin{rem}[近接点アルゴリズム]
非正則化（$\epsilon=0$）問題~\eqref{eq-mk-discr}の解を近似するために、$\KL$計量に対するいわゆる近接点アルゴリズムを用いてSinkhornアルゴリズムを反復的に使用することが可能である。
%
$F(\P) \eqdef \dotp{\P}{\pi} + \iota_{\CouplingsD(\a,\b)}(\P)$を非正則化目的関数とする。
%
$\KLD$ダイバージェンスに対する近接点反復は、$F$の最小値、ひいては非正則化OT問題~\eqref{eq-mk-discr}の解を次の反復計算によって求める
\eql{\label{eq-prox-point}
	\itt{\P} \eqdef \Prox_{\frac{1}{\epsilon} F}^{\KLD}(\it{\P})
	\eqdef
	\uargmin{\P \in \RR_+^{n \times m}} \KLD(\P|\it{\P}) + \frac{1}{\epsilon} F(\P)
}
任意の$\init{\P}$から出発する（~\eqref{eq-prox-kl}も参照）。
%
近接点アルゴリズムは最も基本的な近接分離法である。
%
Euclid計量に対して最初に導入され（例えば(\citealt{rockafellar1976monotone})を参照）、任意のBregman・ダイバージェンスに拡張される~\citep{censor1992proximal}。したがって特に$\KLD$ダイバージェンスに対して適用できる（注意~\ref{rem-bregman}を参照）。
%
近接作用素は通常閉じた形では利用できないため、何らかの形の部分反復が必要である。
%
~\eqref{eq-prox-point}に現れる最適化は、エントロピー正則化問題~\eqref{eq-regularized-discr}と非常に類似しており、負のエントロピー$-\HD$の代わりに相対エントロピー$\KLD(\cdot|\it{\P})$を用いたものである。
%
命題~\ref{prop-regularized-primal}とSinkhorn反復~\eqref{eq-sinkhorn}は、Gibbsカーネルを$\K=e^{-\frac{\C}{\epsilon}} \odot \it{\P} = ( e^{-\frac{\C_{i,j}}{\epsilon}} \it{\P}_{i,j} )_{i,j}$と定義することで、このより一般的な設定にも適用される。
%
したがって反復~\eqref{eq-prox-point}は、各反復でSinkhornアルゴリズムを実行することで実装できる。
%
簡単のため$\init{\P}=\ones_n \ones_m^\top$と仮定すると、これらの反復は次の形になる
\begin{align*}
	\itt{P} &=
		\diag(\it{\uD}) ( e^{-\frac{\C}{\epsilon}} \odot \it{\P} ) \diag(\it{\vD})\\
		&=
		\diag(\it{\uD} \odot \cdots \odot \init{\uD})
			e^{-\frac{(\ell+1)\C}{\epsilon}} \odot \it{\P} )
		\diag(\it{\vD} \odot \cdots \odot \init{\vD}).
\end{align*}
近接点反復はしたがってカーネル$e^{-\frac{\C}{\epsilon/\ell}}$、すなわち減衰する正則化パラメータ$\epsilon/\ell$を用いてSinkhornアルゴリズムを反復的に適用する。
%
この方法は、Sinkhornと正則化の減衰スケジュールを組み合わせた一連の研究と密接に関連している；例えば~\citep{kosowsky1994invisible}を参照。これらは空間次元が小さい場合に、カップリングを適応的なスパースグリッド上で近似するマルチグリッド戦略と組み合わせることで効率的である~\citep{schmitzer2016stabilized}。
\end{rem}
%%%%%%%


%%%%%%%%%
\begin{rem}[他の正則化]
~\eqref{eq-regularized-discr}のエントロピー項$-\HD(\P)$を任意の狭義凸ペナルティ$R(\P)$に置き換えることが可能である。これは例えば~\citep{dessein2016regularized}で詳述されている。典型的な例は$\ell^2$ノルムの2乗
\eql{\label{eq-quad-regul}
	R(\P) = \sum_{i,j} \P_{i,j}^2 + \iota_{\RR_+}(\P_{i,j});
}
である；~\citep{essid2017quadratically}を参照。
%
別の例として、ツァリスエントロピーの族がある~\citep{muzellec2017tsallis}。
%
ただし、ペナルティ関数が$\P$の成分が非正の場合にも定義される場合（例えば2次正則化~\eqref{eq-quad-regul}がそうである）、周辺制約$\P\ones_m=\a$と$\P^\top \ones_n=\b$に加えて非負制約$\P \geq 0$を追加しなければならないことに注意する。実際、エントロピーを用いる場合は非負制約を無視できる。なぜならそのペナルティには$\P$の成分を正の象限に留める対数項が含まれているからである。これは制約集合がもはやアフィンではなくなり、反復Bregman射影が解に収束しなくなることを意味する。
%
回避策として、~\citeauthor{Dykstra83}のアルゴリズム~(\citeyear{Dykstra83,Dykstra85})（\citealt{bauschke-lewis}も参照）を用いることができ、~\citep{2015-benamou-cisc}で詳述されている。このアルゴリズムは$R$に関連するBregman・ダイバージェンスに従った射影を用いる。Bregman・ダイバージェンスの詳細については注意~\ref{rem-bregman}を参照されたい。問題は、一般にこれらの射影を明示的に計算できないことである。
%
2乗ノルム~\eqref{eq-quad-regul}の場合、これは$(\Cc^1_\a,\Cc^2_\b)$へのEuclid射影（追加の正値制約付き）の計算に対応し、単体上の射影アルゴリズム~\citep{condat2015fast}を用いて効率的に解くことができる。2次正則化のエントロピーに対する主な利点は、最適カップリングのスパースな近似を生成することであるが、これはアルゴリズムが遅くなり、複数の最適輸送を同時に計算する際にSinkhornほど効率的に並列化できないという代償を伴う（\S\ref{rem-parallel}で議論）。図~\ref{fig-quad-regul}はエントロピー正則化と2次正則化による近似を対比している。
\end{rem}
%%%%%%%%%


% G B D H
\newcommand{\myfigRegQuad}[1]{\imgBox{\includegraphics[width=.185\linewidth,trim=65 55 45 40,clip]{quadratic-regul/quadratic-levelsets-#1}}}
\newcommand{\myfigRegEntrop}[1]{\imgBox{\includegraphics[width=.185\linewidth,trim=65 55 45 40,clip]{quadratic-regul/entropy-levelsets-#1}}}

\begin{figure}[h!]
\begin{tabular}{@{}c@{}c@{}c@{}c@{}c@{}}
\myfigRegEntrop{1} &
\myfigRegEntrop{3} &
\myfigRegEntrop{4} &
\myfigRegEntrop{7} &
\myfigRegEntrop{9} \\
{\color[rgb]{.6,.6,0} $\varepsilon=10$} &
{\color[rgb]{.45,.6,.15} $\varepsilon=1$} &
{\color[rgb]{.3,.6,.3} $\varepsilon=0.5 \cdot 10^{-1}$} &
{\color[rgb]{.15,.6,.45} $\varepsilon=10^{-1}$} &
{\color[rgb]{0,.6,.6} $\varepsilon=10^{-3}$} \\
\myfigRegQuad{1} &
\myfigRegQuad{3} &
\myfigRegQuad{4} &
\myfigRegQuad{7} &
\myfigRegQuad{9} \\
{\color[rgb]{.6,.6,0} $\varepsilon=5 \cdot 10^3$} &
{\color[rgb]{.45,.6,.15}$\varepsilon=10^3$} &
{\color[rgb]{.3,.6,.3} $\varepsilon=10^2$} &
{\color[rgb]{.15,.6,.45} $\varepsilon=10$} &
{\color[rgb]{0,.6,.6} $\varepsilon=1$}
\end{tabular}
\caption{\label{fig-quad-regul}
エントロピー正則化$R=-\HD$（上段）と2次正則化$R=\norm{\cdot}^2+\iota_{\RR_+}$（下段）の比較。
周辺分布$(\al,\be)$は図~\ref{fig-entropic-copula}と同じ。
}
\end{figure}



%%%%%%%%%%
\begin{rem1}{重心射影}\label{rem-barycenric-proj}
	注意~\ref{rem:entrop-reg}の設定を再び考える。ここでは離散測度間のOTを近似するためにエントロピー正則化を用いる。\eqref{eq-mk-discr}のカントロビッチ定式化とそのエントロピー正則化~\eqref{eq-regularized-discr}はともにカップリング$\P \in \CouplingsD(\a,\b)$を与える。$\Y=\RR^\dim$の場合に輸送写像$\T : \X \rightarrow \Y$を定義するために、いわゆる重心射影写像を定義できる
	\eql{\label{eq-baryproj}
		\T : x_i \in \X \longmapsto \frac{1}{\a_i} \sum_j \P_{i,j} y_j \in \Y,
	}
	ここで入力測度は~\eqref{eq-pair-discr}の形の離散測度である。この写像は$\al$の台の点$(x_i)_i$に対してのみ定義されることに注意する。
	%
	$\T$が置換行列である場合（命題~\ref{prop-matching-kanto}で詳述）、$\T$はMonge写像に等しく、$\varepsilon \rightarrow 0$のとき、重心射影はそれが一意であればその写像に漸近的に収束する。
	%
	任意の（必ずしも離散でない）測度に対して、~\eqref{eq-mk-generic}またはその正則化版~\eqref{eq-entropic-generic}を解くことでカップリング$\pi \in \Couplings(\al,\be)$が定義される。このカップリング$\pi$は常に$\al \otimes \be$に関する密度$\frac{\d\pi(x,y)}{\d\al(x)\d\be(y)}$を持つことに注意する。したがって写像は次の式で復元できる
	\eql{\label{eq-bary-proj}
		\T : x \in \X \longmapsto \int_\Yy y \frac{\d\pi(x,y)}{\d\al(x)\d\be(y)} \d\be(y).
	}
	$\varepsilon=0$の場合に$\pi$がMonge写像のグラフ上に台を持つとき（注意~\ref{rem-exist-mongemap}を参照）、$\varepsilon>0$を用いるとこの写像の滑らかな近似が得られる。
	%
	このような重心射影は、OTMonge写像を画像処理の問題に適用するのに有用である；カラー変換への応用は図~\ref{fig-colors}を参照。また、Wasserstein計量を備えた確率測度の空間における主測地線の近似を計算するためにも使用されている；~\citep{SeguyCuturi}を参照。
\end{rem1}
%%%%%%%%%



%%%%%%%%%%
\begin{rem1}{Hilbert計量}
~\citep{franklin1989scaling}によって最初に説明されたように、Sinkhornの大域的収束解析は$\RR_{+,*}^n$（正ベクトル）上のHilbert射影計量を用いることで大幅に簡略化される。これは次のように定義される
\eq{
	\foralls (\uD,\uD') \in (\RR_{+,*}^n)^2, \quad
	\Hilbert(\uD,\uD') \eqdef \log \umax{i,j} \frac{ \uD_i \uD_{j}' }{ \uD_{j} \uD_{i}'  }.
}
これは射影錐$\RR_{+,*}^n/\sim$上の距離であることが示される。ここで$\uD \sim \uD'$は$\exists r>0, \uD=r\uD'$を意味する（ベクトルはスケーリングの不定性を除いて等しい、ゆえに「射影的」の名がある）。
%
すなわち$\Hilbert$は三角不等式を満たし、$\Hilbert(\uD,\uD')=0$となるのは$\uD \sim \uD'$の場合に限る。
%
これはHilbertの元来の有界開凸集合上の距離~\citep{hilbert1895gerade}の射影版である。
%
射影錐$\RR_{+,*}^n/\sim$はこの距離に関して完備距離空間である。
%
対数変換により、正の錐の半直線上のHilbert計量は変動半ノルム（加法定数の不定性を持つベクトル間のノルム）と等長である
\eql{\label{eq-hilbert-var}
	\Hilbert(\uD,\uD') = \norm{\log(\uD)-\log(\uD')}_{\text{var}}
}
\eq{
	\qwhereq
	\norm{\fD}_{\text{var}} \eqdef (\max_i \fD_i) - (\min_i \fD_i).
}
%
この変動半ノルムは$\ell^\infty$ノルムと密接に関連しており、常に$\norm{\fD}_{\text{var}}  \leq 2 \norm{\fD}_\infty$が成り立つ。ある固定された$i$に対して$\fD_i=0$を課すと、$\norm{\fD}_\infty \leq \norm{\fD}_{\text{var}}$であるから逆向きの不等式も成り立つ。これらの評価はSinkhornの収束解析に特に有用である（以下の注意~\ref{rem-global-conv-sinkh}を参照）。なぜなら~\eqref{eq-dualsinkhorn-constraints2}を解く双対変数$\fD=\log(\uD)$は加法定数の不定性を持つため、ある$i$に対して$\fD_i=0$を課すことができるからである。
%
Hilbert計量は~\citep{birkhoff1957extensions}と~\citep{samelson1957perron}によって独立に導入された。彼らは以下の基本定理を証明した。これは正行列が正ベクトルの錐上で厳密な縮小写像であることを示している。


\begin{thm}\label{thm-birkoff}
	$\K \in \RR_{+,*}^{n \times m}$とする；すると$(\vD,\vD') \in (\RR_{+,*}^m)^2$に対して
	\eq{
		\Hilbert(\K \vD,\K \vD') \leq \la(\K) \Hilbert(\vD,\vD'),
		\text{ where }
		\choice{
			\la(\K) \eqdef \frac{ \sqrt{\eta(\K)}-1 }{ \sqrt{\eta(\K)}+1 } < 1, \\
			\eta(\K) \eqdef \umax{i,j,k,\ell} \frac{ \K_{i,k} \K_{j,\ell} }{ \K_{j,k} \K_{i,\ell} }.
		}
	}
\end{thm}

図~\ref{fig-h-m}はこの定理を図解している。
\end{rem1}


\begin{figure}[h!]
\centering
\includegraphics[width=.22\linewidth]{hilbert-metric/hilbert-metric} \qquad
\includegraphics[width=.7\linewidth]{hilbert-metric/cones}\\
\caption{\label{fig-h-m}
左：Hilbert計量$\Hilbert$は錐内の半直線上の距離である（ここでは正ベクトル）。
右：正行列$\K$の反復によって誘導される縮小の可視化。
}
\end{figure}


%%%%%%%%%%
\begin{rem1}{Perron-Frobenius}
定理~\ref{thm-birkoff}の典型的な応用は、Perron-Frobeniusの定理の定量的な証明を与えることである。注意~\ref{rem-local-conv}で説明するように、これはSinkhorn反復の局所的な線形化に関連している。$\K^\top \ones_n = \ones_n$を満たす行列$\K \in \RR_+^{n \times n}$は$\simplex_n$を$\simplex_n$に写す。さらに$\K>0$であれば、定理~\ref{thm-birkoff}により計量$\Hilbert$に関して厳密に縮小的であるから、$\K p^\star=p^\star$を満たす一意の不変確率分布$p^\star \in \simplex_n$が存在する。さらに、任意の$p_0 \in \simplex_n$に対して$\Hilbert(\K^\ell p_0,p^\star) \leq \la(\K)^\ell \Hilbert(p_0,p^\star)$、すなわち行列の反復が$p^\star$に線形収束する。これは図~\ref{fig-perron}に図示されている。
\end{rem1}
%%%%%%%%%%

\begin{figure}[h!]
\centering
\begin{tabular}{@{}c@{}c@{}c@{\hspace{10mm}}c@{}}
\includegraphics[width=.22\linewidth]{hilbert-metric/simplex-0}&
\includegraphics[width=.22\linewidth]{hilbert-metric/simplex-1}&
\includegraphics[width=.22\linewidth]{hilbert-metric/simplex-2}&
\includegraphics[width=.22\linewidth]{hilbert-metric/simplex-evol}\\
$\simplex_3$ & $\K \simplex_3$ & $\K^2 \simplex_3$ &  $\{ \K^\ell \simplex_3 \}_\ell$
\end{tabular}
\caption{\label{fig-perron}
$\K^\top \ones_3=\ones_3$を満たす$\K \in \RR_{+,*}^{3 \times 3}$の不変確率分布$p^\star$への$\K^\ell \simplex_3 \rightarrow \{p^\star\}$の変化。
}
\end{figure}




%%%%%%%%%%
\begin{rem11}{大域的収束}{rem-global-conv-sinkh}
次の定理は~\citep{franklin1989scaling}によって証明され、定理~\ref{thm-birkoff}を用いてSinkhorn反復の線形収束を示す。

\begin{thm}
	$(\it{\uD},\it{\vD}) \rightarrow (\uD^\star,\vD^\star)$であり
	\eql{\label{eq-convlin-sinkh}
		\Hilbert(\it{\uD}, \uD^\star) = O(\la(\K)^{2\ell}), \quad
		\Hilbert(\it{\vD}, \vD^\star) = O(\la(\K)^{2\ell}).
	}
	また
	\eql{\label{eq-convsinkh-control}
		\begin{split}
		\Hilbert(\it{\uD}, \uD^\star) &\leq \frac{\Hilbert( \it{\P}\ones_m,\a )}{1-\la(\K)^2}, \\
		\Hilbert(\it{\vD}, \vD^\star) &\leq \frac{\Hilbert( \P^{(\ell),\top} \ones_n,\b )}{1-\la(\K)^2},
		\end{split}
	}
	ここで$\it{\P} \eqdef \diag(\it{\uD}) \K \diag(\it{\vD})$と表した。最後に
	\eql{\label{eq-convlin-sinkh-prim}
		\|\log(\it{\P}) - \log(\P^\star)\|_\infty \leq \Hilbert(\it{\uD}, \uD^\star) + \Hilbert(\it{\vD}, \vD^\star),
	}
	ここで$\P^\star$は~\eqref{eq-regularized-discr}の一意の解である。
\end{thm}

\begin{proof}
	任意の$(\vD,\vD') \in (\RR_{+,*}^m)^2$に対して
	\eq{
		\Hilbert(\vD,\vD') = \Hilbert(\vD/\vD',\ones_m) = \Hilbert(\ones_m/\vD,\ones_m/\vD').
	}
	であることに注意する。これにより
	\begin{align*}
		\Hilbert(\itt{\uD},\uD^\star) &= \Hilbert\pa{ \frac{\a}{\K \it{\vD}}, \frac{\a}{\K \vD^\star} } \\
		&= \Hilbert( \K \it{\vD}, \K \vD^\star ) \leq \la(\K) \Hilbert( \it{\vD}, \vD^\star ),
	\end{align*}
	ここで定理~\ref{thm-birkoff}を用いた。これにより~\eqref{eq-convlin-sinkh}が示される。また三角不等式を用いると
	\begin{align*}
		\Hilbert(\it{\uD},\uD^\star) &\leq \Hilbert(\itt{\uD},\it{\uD}) + \Hilbert(\itt{\uD},\uD^\star) \\
		&\leq \Hilbert\pa{ \frac{\a}{\K \it{\vD}},\it{\uD} } + \la(\K)^2 \Hilbert(\it{\uD},\uD^\star) \\
		&= \Hilbert\pa{ \a,\it{\uD} \odot  ( \K \it{\vD} ) } + \la(\K)^2 \Hilbert(\it{\uD},\uD^\star),
	\end{align*}
	$\it{\uD} \odot  ( \K \it{\vD} ) = \it{\P}\ones_m$であるから、~\eqref{eq-convsinkh-control}の第一部分が得られる（第二部分も同様）。
	%
	~\eqref{eq-convlin-sinkh-prim}の証明は~\citep[Lem. 3]{franklin1989scaling}に従う。
\end{proof}

評価~\eqref{eq-convsinkh-control}は、周辺制約違反に関するある種の誤差測度、例えば$\| \it{\P} \ones_m - \a \|_1$と$\|\transp{\it{\P}} \ones_n - \b \|_1$が、収束を監視するための有用な停止基準であることを示している。
%
~\eqref{eq-hilbert-var}のおかげで、スケーリング変数$(\it{\uD},\it{\vD})$に関するこれらのHilbert計量の収束率は、双対変数$(\it{\fD},\it{\gD}) \eqdef (\varepsilon\log(\it{\uD}),\varepsilon\log(\it{\vD}))$に対する変動ノルム$\norm{\cdot}_{\text{var}}$での線形収束率を与える。

図~\ref{fig-sinkhorn-convergence}の下段は、制約違反に関するこの線形収束率を強調し、$\varepsilon\rightarrow 0$でこの収束率がどのように悪化するかを示している。
%
これらの結果は~\citep{franklin1989scaling}で証明されており、非線形Perron-Frobenius理論~\citep{lemmens2012nonlinear}と密接に関連している。Perron-Frobenius理論は反復の線形化に対応する；~\eqref{eq-linearized-sinkh}を参照。この収束解析は~\citep{linial1998deterministic}によって拡張され、Sinkhornの各反復がスケーリングされたカップリング行列のパーマネントを増加させることが示された。
\end{rem11}
%%%%%%%%%%



%%%%%%%%%%
\begin{rem1}{局所収束}\label{rem-local-conv}
大域的線形収束率~\eqref{eq-convlin-sinkh-prim}はしばしばかなり悲観的であり、特に$\varepsilon=0$のときにMonge写像が存在する場合の$\X=\Y=\RR^d$においてそうである（注意~\ref{rem-monge2}を参照）。大域的収束率は、コスト行列$\C$がランダムに近いより困難な状況ではむしろ鋭く、これらの場合、収束率は$\varepsilon$に関して指数的に悪化し、$1-\la(\K) \sim e^{-1/\varepsilon}$となる。
%
より精密な漸近的収束解析を得るために（例えば高精度の解に関心があり、多数の反復を行う場合）、通常は局所収束率を研究する。
%
Sinkhornの更新を不動点写像$\itt{\fD} = \Phi( \it{\fD} )$の反復として書くことができる。ここで
\eq{
	\Phi \eqdef \Phi_2 \odot \Phi_1
	\qwhereq
	\choice{
		\Phi_1(\fD) = \varepsilon \log \transp{\K} ( e^{\fD/\varepsilon} ) - \log(\b), \\
		\Phi_2(\gD) = \varepsilon \log \K ( e^{\gD/\varepsilon} ) - \log(\a).
	}
}
~\eqref{eq-dual-formulation}を解く最適な$(\fD,\gD)$に対して、$\P = \diag(e^{\fD/\varepsilon}) \K \diag(e^{\gD/\varepsilon})$を~\eqref{eq-regularized-discr}を解く最適カップリングとすると、次のJacobianが得られる：
\eql{\label{eq-linearized-sinkh}
	\partial \Phi(\fD) = \diag(\a)^{-1} \odot \P \odot \diag(\b)^{-1} \odot \transp{\P}.
}
このJacobianは$\partial \Phi(\fD) \ones_n = \ones_n$を満たす正行列であり、したがってPerron-Frobeniusの定理により、固有値$1$に対応する支配的な固有ベクトル$\ones_m$を1つ持つ。$\fD$は定数の不定性を持つため、実際には第2固有値$1-\kappa<1$が局所的な線形収束率を支配し、十分大きな$\ell$に対して
\eq{
	\|\it{\fD}-\fD \| = O( (1-\kappa)^\ell ).
}
%
数値的には、「単純な場合」（例えば$\varepsilon=0$のとき滑らかなMonge写像が存在する場合）において、この収束率は$\kappa \sim \varepsilon$のようにスケーリングする。
%
二重確率的（割当）場合の詳細については~\citep{knight2008sinkhorn}を参照されたい。
\end{rem1}
%%%%%%%%%%






%%%%%%%%%%%%%%%%%%%%%%%%%%%%%%%%%%%%%%%%%%%%%%%%%%%%%%%%%%%%%%%%%%%%%%%%%%%%%
\section{Sinkhorn反復の高速化}

Sinkhorn反復の主な計算ボトルネックは、カーネル$\K$と$\K^\top$に対するベクトル-行列積であり、単純に実装すると計算量は$O(nm)$である。以下では、計算量を大幅に改善できるいくつかの重要なケースを詳述する。

%%%%%%%
\begin{rem}[並列かつGPU向きの計算]\label{rem-parallel} Sinkhornアルゴリズムの簡潔さは、ヒストグラムのペア間の複数の正則化Wasserstein距離を同時に計算するための極めて効率的なアプローチを提供する。
$N$を整数とし、$\a_1,\ldots,\a_N$を$\simplex_n$のヒストグラム、$\b_1,\ldots,\b_N$を$\Sigma_m$のヒストグラムとする。
すべての$N$個の近似距離$\MKD_\C^\varepsilon(\a_1,\b_1),\dots,\MKD_\C^\varepsilon(\a_N,\b_N)$を計算したい。
この場合、すべてのヒストグラムを格納する$n\times N$と$m\times N$の行列を$\A=[\a_1,\dots,\a_N]$と$\B=[\b_1,\dots,\b_N]$と書くと、これら$N$ペアに対するすべてのSinkhorn反復は、例えば次のように設定することで並列に実行できる
\eql{\label{eq-sinkhorn-par}
	\itt{\UD} \eqdef \frac{\A}{\K \it{\VD}}
	\qandq
	\itt{\VD} \eqdef \frac{\B}{\transp{\K}\itt{\UD}},
}
初期値は$\init{\VD} = \ones_{m \times N}$である。ここで$\frac{\cdot}{\cdot}$は行列の成分ごとの除算に対応する。
$\VD$と$\UD$の収束後、正則化距離の（行）ベクトルは次のように簡略化されることが確認できる
$$ \transp{\ones_n} (    \UD\odot \log \UD \odot (   (\K\odot \C) \VD) +  \UD\odot(   (\K\odot \C)(\VD\odot \log \VD)    ) ) \in\RR^{N}.$$
式~\eqref{eq-sinkhorn}で記述される基本的なSinkhorn反復は本質的にGPU向きであることに注意する。なぜならそれらは行列-ベクトル積のみから構成されるからであり、これは例えば~\citet{slomp2011gpu}でマッチング問題を解くために利用された。しかし、式~\eqref{eq-sinkhorn-par}で示される行列-行列演算はさらに良い並列化の機会を提供し、これが大規模なヒストグラム間のOT距離を計算する際のSinkhornアルゴリズムの成功を説明している。
\end{rem}
%%%%%%%

%%%%%%%%%
\begin{rem}[分離可能カーネルに対する高速化]\label{rem-separable} 本節では、各Sinkhorn反復の計算量を大幅に削減できる重要な特殊ケースを考える。このケースは、コスト行列で考慮される各添字$i$と$j$が$d$個の有限集合$\range{n_1},\dots,\range{n_d}$の直積から取られた$d$-組として記述できる場合に発生する
$$i=(i_k)_{k=1}^d, j=(j_k)_{k=1}^d \in \range{n_1}\times\dots\times\range{n_d}.$$
この設定において、添字$i$と$j$の間のコスト$\C_{ij}$がこれらの部分添字に沿って加法的である場合、すなわちそれぞれサイズ$n_1\times n_,\dots,n_d\times n_d$の$d$個の行列$\C^1,\dots,\C^d$が存在して
$$\C_{ij} = \sum_{k=1}^d \C^k_{i_k,j_k},$$
となる場合、直接的な帰結として、Sinkhorn反復に現れるカーネルは分離可能な乗法構造を持つ
\eql{\label{eq-speedup-separable}
\K_{i,j} = \prod_{k=1}^d \K^k_{i_k,j_k}.
}
このような分離可能な乗法構造により、$\K \uD$の非常に高速な（厳密な）評価が可能になる。実際、$\K$をサイズ$n\times n$の行列として実体化する代わりに（$n=\prod_k n_k$は通常次元$d$に関して指数的であるため禁止的なサイズとなる）、各「スライス」に沿って$\uD$に$\K^k$を適用するだけで$\K\uD$を復元できる。$n=m$の場合、計算量は$O(n^2)$の代わりに$O(n^{1+1/d})$に減少する。
%

この高速化の重要な例は、$\Xx=\Yy=[0,1]^d$で基底コストが$q$-ノルムの$q$乗、$$\c(x,y)=\norm{x-y}^q_q=\sum_{i=1}^d |x_i-y_i|^q, \; q>0;$$であり、空間が正則格子で離散化されて$x_i = (i_1/n_1,\ldots,i_d/n_d)$（$i=(i_1,\dots,i_d)\in\range{n_1}\times \dots\times\range{n_d}$）の点のみが考慮される場合に生じる。この場合、$\K$との積は各1次元$n_k\times n_k$畳み込み行列$$\K^k = \begin{bmatrix}\exp(-\left|\frac{r-s}{n_k}\right|^q/ \varepsilon)\end{bmatrix}_{1 \leq r,s\leq n_k}$$を、$k$番目の添字集合に一致するように第1次元を置換したテンソルとして再形成した$\uD$に適用することで、より効率的に実行できる。
%
例えば、$d=2$（平面の場合）かつ$q=2$（2-Wasserstein、Gauss畳み込みに帰着）の場合、ヒストグラム$\a$、したがってSinkhorn乗数$\uD$は$n_1\times n_2$の行列として表現できる。乗数$\uD$が長さ$n_1n_2$のベクトルではなく$n_1\times n_2$の行列として再形成されていることを強調するために$\mathbf{U}$と書く。すると、単純な実装では$(n_1 n_2)^2$演算を必要とする$\K\uD$の計算は、2つの1次元畳み込みを個別に適用して$$(\K^2(\K^1 \mathbf{U})^T)^T=\K^1\mathbf{U}\K^2,$$として得ることができ、$n_1^2n_2^2$演算の代わりに$(n_1^2)n_2+n_1(n_2^2)$演算で$n_1\times n_2$行列が復元される。この例は上で与えた指数$(1+1/d)$と一致することに注意する。$d$がより大きい場合、（適切なサイズのテンソルとして再形成された）$\uD$の各スライスにこれと全く同じ1次元畳み込みを適用する必要があるが、この演算はGPU上で極めて効率的である。
%

この重要な観察は、2次元や3次元の形状データに最適輸送を適用する際の多くの実用的な成功の基盤となっている。~\citep{2015-solomon-siggraph,2016-bonneel-barycoord}で強調されているように、$200^3=8 \times 10^6$もの大きさの格子上に台を持つ分布が扱われている。
\end{rem}
%%%%%%%%%

%%%%%%%%%
\begin{rem}[近似畳み込み]\label{rem-convol-sinkh}
Sinkhorn反復~\eqref{eq-sinkhorn}の主な計算ボトルネックは、$\K$またはその随伴によるベクトルの乗算にある。分離性~\eqref{eq-speedup-separable}を利用する以外にも、カーネルの他の特殊構造を利用することが可能である。最も単純なケースは平行移動不変カーネル$\K_{i,j} = k_{i-j}$であり、これはEuclid空間$\X=\RR^\dim$の固定された一様格子上で測度を離散化する場合に典型的である。このとき$\K \vD = k \star \vD$は畳み込みであり、ほぼ線形時間で畳み込みを近似するいくつかのアルゴリズムが存在する。
%
最も一般的なのはFourier変換$\Ff$による方法であり、簡単のため周期境界条件を仮定すると、$\Ff(k \star \vD)=\Ff(k) \odot \Ff(\vD)$であるからである。しかしこれは不安定な計算につながり、小さな$\varepsilon$に対してはしばしば許容できない。
%
計算を高速化するもう1つの一般的な方法は、例えばデリシュフィルタリング法~\citet{deriche1993recursively}を用いて、自己回帰フィルターの連続適用により畳み込みを近似することである。様々な高速フィルタリング法の比較については~\citep{getreuer2013survey}を参照されたい。
\end{rem}
%%%%%%%%%

%%%%%%%%%
\begin{rem}[熱近似における測地線]\label{rem-geod-heat}
非平面領域では、カーネル$\K$は畳み込みではないが、コストが$\C_{i,j}=d_\Mm(x_i,y_j)^p$（$d_\Mm$は曲面$\Mm$（またはより一般的な多様体）上の測地線距離）の場合、$\K=e^{-\frac{d_\Mm}{\varepsilon}}$のベクトルへの適用の高速近似を行うことも可能である。実際、\citeauthor{varadhan-1967}の公式~\citeyearpar{varadhan-1967}は、このカーネルがLaplacianカーネル（$p=1$の場合）および熱核（$p=2$の場合）に近いことを主張している。
%
ヴァラダンの第1公式は次のように述べる
\eql{\label{eq-varadhan-1}
	-\frac{\sqrt{t}}{2} \log( \Pp_t(x,y) ) = d_\Mm(x,y) + o(t)
	\qwhereq
	\Pp_t \eqdef (\Id - t\De_{\Mm})^{-1},
}
ここで$\De_{\Mm}$は多様体$\Mm$に関連するラプラス-ベルトラミ作用素（半負定値）であり、
$\Pp_t$は積分核であり、$g = \int_\Mm \Pp_t(x,y) f(y) \d y$は$g - t \De_{\Mm} g = f$の解である。
%
ヴァラダンの第2公式は次のように述べる
\eql{\label{eq-varadhan-2}
	\sqrt{ -4 t \log( \Hh_t(x,y) ) } = d_\Mm(x,y) + o(t),
}
ここで$\Hh_t$は、$g_t = \int_\Mm \Hh_t(x,y) f(y) \d y$が時刻$t$における熱方程式
\eq{
	\frac{\partial g_t(x)}{\partial t} = (\Delta_\Mm g_t)(x).
}
の解であるように定義された積分核である。
公式~\eqref{eq-varadhan-1}と~\eqref{eq-varadhan-2}の収束はコンパクト多様体上で一様である。
%
数値的には、領域$\Mm$は離散化され（例えば有限要素法を用いて）、$\De_\Mm$は離散Laplacian行列$L$で近似される。典型的な例は区分線形有限要素を用いる場合であり、$L$は有名なコタンジェントLaplacianとなる（この構成の詳細については~\citep{botsch-2010}を参照）。
%
これらの公式はGibbsカーネル$\K_{i,j}=e^{-\frac{d(x_i,y_j)^p}{\varepsilon}}$との積を効率的に近似するために使用できる。
%
式~\eqref{eq-varadhan-1}は、$p=1$の場合に$\varepsilon=\frac{\sqrt{t}}{2}$とし、$\K$との積を$(\Id-t L)^{-1}$との積で置き換えることを示唆する。これには正定値対称線形系の解法が必要である。
%
式~\eqref{eq-varadhan-2}は、熱流の安定的な解法のための陰的Euler法の$R$ステップと組み合わせると、$p=2$の場合に$\K$との積を$4 t=\varepsilon$として$(\Id-\frac{t}{R} L)^{-R}$との積で置き換えることを示唆する。これには$R$回の線形系の解法が必要である。幸い、これらの線形系は各Sinkhorn反復で解かれるため、スパースなコレスキー分解を事前に計算することで効率的に解ける。行列の行と列の並べ替え~\citep{george1989evolution}を行うことで、2次元多様体に対してほぼ線形のスパース性が得られ、各Sinkhorn反復は線形計算量を持つ（性能は多様体の次元とともに低下する）。
%
ヴァラダンの公式を測地線距離の近似に使用することは~\citep{Crane2013}で最初に提案され、Sinkhorn反復との併用は~\citep{2015-solomon-siggraph}で提案された。
%
\todoK{Maybe add some pedagological figure illustrating the method.}
\end{rem}
%%%%%%%%%


%%%%%%%%%%
\begin{rem}[外挿加速]
	Sinkhornアルゴリズムは（注意~\ref{rem-local-conv}で示されたように）不動点アルゴリズムであるため、標準的な線形または非線形外挿スキームを用いて解の近傍での不動点写像の条件数を改善し、線形収束率を向上させることができる。
	%
	これは逐次過緩和法（例えば~\citep{hadjidimos2000successive}を参照）と類似しており、局所的な線形収束率はある$\kappa>0$に対して$O((1-\kappa)^\ell)$から$O((1-\sqrt{\kappa})^\ell)$に改善される（注意~\ref{rem-local-conv}を参照）。
	%
	詳細については~\citep{2016-peyre-qot}を参照されたい\todoK{citer papier recent Lenaic}。
\end{rem}
%%%%%%%%%%



%%%%%%%%%%%%%%%%%%%%%%%%%%%%%%%%%%%%%%%%%%%%
\section{安定性と対数領域での計算}

注意~\ref{rem-stability}で簡単に述べたように、Sinkhornアルゴリズムは正則化パラメータ$\epsilon$がコスト行列$\C$の成分に比べて小さい場合に数値的なオーバーフローを起こす。この懸念は計算を対数領域で行うことである程度軽減できる。このアプローチの妥当性は、~\eqref{eq-regularized-discr}に関連する双対問題を考えることでより明確になる。この双対問題では対数領域での計算が自然に現れる。

\begin{prop}
次が成り立つ
\eql{\label{eq-dual-formulation}
	\MKD_\C^\varepsilon(\a,\b) = \umax{\fD \in \RR^n,\gD \in \RR^m}
		 \dotp{\fD}{\a} + \dotp{\gD}{\b}
		% - \varepsilon\sum_{i,j} e^{ \frac{\fD_i+\gD_j-\C_{i,j}}{\varepsilon} }.
		- \varepsilon \dotp{e^{\fD/\varepsilon} }{ \K e^{\gD/\varepsilon}}.
}
%
最適な$(\fD,\gD)$は~\eqref{eq-scaling-form}に現れるスケーリング$(\uD,\vD)$と次のように関連する
\eql{\label{eq-entropy-pd}
	(\uD,\vD)=(e^{\fD/\varepsilon},e^{\gD/\varepsilon}).
}
\end{prop}

\begin{proof}
命題~\ref{prop-regularized-primal}の証明の終わりから出発する。そこでは最適な主問題の解$\P$と周辺制約に対する双対乗数$\fD$と$\gD$が次のように関連づけられる
\eq{
	\label{eq-prim-dual-ent}\P_{i,j}=e^{\fD_i/\varepsilon}e^{-\C_{i,j}/\varepsilon}e^{\gD_j/\varepsilon}.
}
式~\eqref{eq-sinkhorn-lagrangian}のLagrangian$\Lag(\P,\fD,\gD)$に、$\fD$と$\gD$の関数としての最適な$\P$を代入すると、Lagrange双対関数は次に等しくなる
\eql{\label{eq-lagrange-dual-naive}
	\fD,\gD \mapsto \dotp{e^{\fD/\varepsilon}}{\left(\K\odot \C\right)e^{\gD/\varepsilon}} - \varepsilon \HD(\diag(e^{\fD/\varepsilon}) \K \diag(e^{\gD/\varepsilon})).
}
$\varepsilon$でスケーリングした$\P$の負エントロピー、すなわち$\varepsilon \dotp{\P}{\log \P - \ones_{n\times m}}$は、$\fD,\gD, \C$の関数として明示的に表すことができ
\begin{align*}
&\dotp{\diag(e^{\fD/\varepsilon}) \K \diag(e^{\gD/\varepsilon})}{\fD\transp{\ones_m}+\ones_n\transp{\gD}-\C-\varepsilon\ones_{n\times m}}\\
&= -\dotp{e^{\fD/\varepsilon}}{\left(\K\odot \C\right)e^{\gD/\varepsilon}} + \dotp{\fD}{\a}+ \dotp{\gD}{\b}-\varepsilon \dotp{e^{\fD/\varepsilon}}{\K e^{\gD/\varepsilon}};
\end{align*}
したがって~\eqref{eq-lagrange-dual-naive}の第1項は上のエントロピーの第1項と相殺される。残りの項が~\eqref{eq-dual-formulation}に現れるものである。
\end{proof}



%%%%%%%%%%
\begin{rem}[双対問題に対するブロック座標上昇としてのSinkhorn]\label{rem-sinkh-block}
制約なし最大化問題~\eqref{eq-dual-formulation}を解く単純なアプローチは、正確な\emph{ブロック座標上昇}戦略を用いることである。すなわち$\fD$と$\gD$を交互に更新して、~\eqref{eq-dual-formulation}の目的関数のこれらの変数に関するそれぞれの勾配をゼロにする。実際、いくつかの基本的な計算の後、~\eqref{eq-dual-formulation}の目的関数を$Q(\fD,\gD)$と書くと
\begin{align}
	\label{eq-dualupadate-sinkh-1}\nabla|_\fD\, Q(\fD,\gD) &=  \a - e^{\fD/\varepsilon}\odot \left(\K e^{\gD/\varepsilon}\right),\\
	\label{eq-dualupadate-sinkh-2}\nabla|_\gD\, Q(\fD,\gD) &=  \b - e^{\gD/\varepsilon}\odot \left(\transp{\K} e^{\fD/\varepsilon}\right).
\end{align}
であることがわかる。ブロック座標上昇はしたがって、任意の$\init{\gD}$から出発して、$l\geq 0$に対して以下の更新を逐次適用することにより閉じた形で実装できる：
\begin{align}
	\label{eq-slse-sinkh-1}\itt{\fD} &= \varepsilon \log \a -\varepsilon\log\left(\K e^{\it{\gD}/\varepsilon}\right), \\
	\label{eq-slse-sinkh-2}\itt{\gD} &= \varepsilon \log \b- \varepsilon\log\left(\transp{\K} e^{\itt{\fD}/\varepsilon}\right).
\end{align}
このような反復は、~\eqref{eq-entropy-pd}で強調された主-双対関係を考慮すると、Sinkhorn反復~\eqref{eq-sinkhorn}と数学的に等価である。実際、任意の反復において
\eq{
	(\it{\fD},\it{\gD}) = \varepsilon ( \log(\it{\uD}), \log(\it{\vD}) ).
}
が回復される。
\end{rem}
%%%%%%%%%%

%%%%%%%%%%
\begin{rem}[ソフト最小値書き換え]
反復~\eqref{eq-slse-sinkh-1}と~\eqref{eq-slse-sinkh-2}は、以下の表記法を用いて別の解釈を与えることができる。実数のベクトル$\z$が与えられたとき、その座標の\emph{ソフト最小値}を$\smine\z$と書く。すなわち
\eql{\label{eq-defn-softmin}
	\smine\z= -\varepsilon \log \sum_i e^{-\z_i/\varepsilon}.
}
$\varepsilon\rightarrow 0$のとき、任意のベクトル$\z$に対して$\smine(\z)$は$\min\,\z$に収束することに注意する。実際、$\smine$は図~\ref{fig-softmin}に示すように、$\min$関数の微分可能な近似として解釈できる。

\begin{figure}[h!]
\centering
\begin{tabular}{@{}c@{\hspace{1mm}}c@{\hspace{1mm}}c@{\hspace{1mm}}c@{\hspace{1mm}}c@{}}

\includegraphics[width=.19\linewidth]{soft-min/lse-1}&
\includegraphics[width=.19\linewidth]{soft-min/lse-05}&
\includegraphics[width=.19\linewidth]{soft-min/lse-01}&
\includegraphics[width=.19\linewidth]{soft-min/lse-001}&
\includegraphics[width=.19\linewidth]{soft-min/lse-0001} \\
$\varepsilon=1$ &
$\varepsilon=0.5$ &
$\varepsilon=10^{-1}$ &
$\varepsilon=10^{-2}$ &
$\varepsilon=10^{-3}$
\end{tabular}
\caption{\label{fig-softmin}
2次元における関数$\smine(\z)$の表示、$\z \in \RR^2$、様々な$\varepsilon$に対して。
}
\end{figure}

この表記法を用いると、式~\eqref{eq-slse-sinkh-1}と~\eqref{eq-slse-sinkh-2}は次のように書き換えられる
\begin{align}
\label{eq-slse-smin-1}(\itt{\fD})_i &= \smine\, (\C_{ij}-\it{\gD}_j)_j + \varepsilon \log \a_i, \\
\label{eq-slse-smin-2}(\itt{\gD})_j &= \smine\, (\C_{ij}-\it{\fD}_i)_i + \varepsilon \log \b_j.
\end{align}
ここで$\smine\, (\C_{ij}-\it{\gD}_j)_j$は行列$(\C-\ones_n (\it{\gD})^\top )$の$j$列のすべての値のソフト最小値を表す。表記を簡略化するために、行列を入力として取り、その列または行のソフト最小値の列ベクトルを出力する演算子を導入する。すなわち、任意の行列$A\in\RR^{n\times m}$に対して
\begin{align*}
	\SMINEr(\A) \eqdef \pa{\smine \pa{\A_{i,j}}_j}_i\in\RR^n,\\
	\SMINEc(\A) \eqdef \pa{\smine \pa{\A_{i,j}}_i}_j\in\RR^m.
\end{align*}
と定義する。これらの演算は\S\ref{sec-semi-discr-entropy}で導入されるエントロピー的$\c$-変換と等価であることに注意する（特に~\eqref{eq-sinkh-c-transf}を参照）。
%
この表記法を用いると、Sinkhornの反復は次のように書ける
\begin{align}
	\label{eq-lse-dual-1}\itt{\fD} &= \SMINEr\, (\C-\ones_n\transp{\it{\gD}}) + \varepsilon \log \a, \\
	\label{eq-lse-dual-2}\itt{\gD} &= \SMINEc\, (\C-\it{\fD}\transp{\ones_m})  + \varepsilon \log \b.
\end{align}
$\varepsilon \rightarrow 0$のとき$\smine$は$\min$に収束するが、$\varepsilon=0$の極限では反復はもはや収束しない。なぜなら交互最小化は制約付き問題では収束せず、これは非正則化双対~\eqref{eq-dual}の場合だからである。
\end{rem}
%%%%%%%%%%



%%%%%%%%%%
\begin{rem}[対数領域Sinkhorn]\label{rem-log-sinkh}
Sinkhorn更新~\eqref{eq-sinkhorn}と数学的に等価であるが、反復~\eqref{eq-slse-smin-1}と~\eqref{eq-slse-smin-2}は、小さな$\varepsilon$の値に対するアンダーフローを回避するために\emph{ログサムエクスポネンシャル}安定化トリックの使用を示唆する。$\underbar{z}=\min \z$と書くと、そのトリックは$\smine \z$を次のように評価することを示唆する
\eql{\label{eq-log-domain-min}
	\smine \z = \underbar{z} -\varepsilon \log \sum_i e^{-(\z_i-\underbar{z})/\varepsilon}.
}
~\eqref{eq-log-domain-min}のように$\underbar{z}$を引いて対数領域反復を安定化する代わりに、実際には以前に計算されたスケーリングを引くことができる。
これは安定化された反復につながる
\begin{align}
	\label{eq-lse-sinkh-1}\itt{\fD} &= \SMINEr( \logP(\it{\fD},\it{\gD}))  + \it{\fD} + \varepsilon\log(\a), \\
	\label{eq-lse-sinkh-2}\itt{\gD} &= \SMINEc( \logP(\itt{\fD},\it{\gD})) + \it{\gD} + \varepsilon \log(\b),
\end{align}
ここで
\eq{
	\logP(\fD,\gD) = \pa{\C_{i,j} - \fD_i - \gD_j}_{i,j}.
}
と定義した。
% The log-sum-exp operators are defined as
% \begin{align*}
% 	\foralls i \in \range{n}, \quad \LSE_1(\logP)_i &= \varepsilon \log \sum_j \exp( \logP_{i,j} ), \\
% 	\foralls j \in \range{m}, \quad \LSE_2(\logP)_j &= \varepsilon \log \sum_i \exp( \logP_{i,j} ).
% \end{align*}
もとの反復~\eqref{eq-sinkhorn}とは対照的に、これらの対数領域反復~\eqref{eq-lse-sinkh-1}と~\eqref{eq-lse-sinkh-2}は任意の$\varepsilon>0$に対して安定である。
なぜなら量$\logP(\fD,\gD)$は反復中に有界に留まるからである。
欠点は、各ステップで$nm$回の$\exp$の計算が必要なことである。
$\SMINEr$または$\SMINEc$の計算は、行列積よりも通常大幅に遅く、行列$\logP$の行ごとのソフト最小値の計算を必要とする。したがって、複数の周辺分布に対するSinkhorn写像の適用を効率的に並列化する方法は存在しない。
%
小さな次元のEuclid領域では、計算を大幅に高速化するために、減衰する$\varepsilon$戦略とスパースグリッドを用いた効率的なマルチスケール解法を開発することが可能である~\citep{schmitzer2016stabilized}。
\end{rem}
%%%%%%%%%%




%%%%%%%%%%
\begin{rem2}{一般的な測度に対する双対}
一般的な（必ずしも離散でない）入力測度$(\al,\be)$に対して、双対問題~\eqref{eq-dual-formulation}は次のようになる
\eql{\label{eq-dual-entropic}
	\usup{(\f,\g) \in \Cc(\X)\times\Cc(\Y)} \int_\X \f \d\al + \int_\Y \g \d\be
		 - \varepsilon \int_{\X\times\Y} e^{ \frac{-c(x,y)+f(x)+g(y)}{\varepsilon} } \d\al(x)\d\be(y).
}
これはもとの問題~\eqref{eq-dual-generic}に現れる制約$\Potentials(\c)$の平滑化に対応し、$\varepsilon \rightarrow 0$の極限で復元される。
%
エントロピー輸送の場合にこれらのカントロビッチポテンシャル$(\f,\g)$の存在（すなわち上限が実際に最大値であること）を証明することは、古典的なOTの場合より容易ではない。なぜなら$c$-変換を使うことができず、ポテンシャルは自動的にLipschitzにならないからである。存在の証明はSinkhorn反復の収束を用いて行うことができる；詳細は~\citep{2016-chizat-sinkhorn}を参照。
\end{rem2}
%%%%%%%%%%




%%%%%%%
\begin{rem2}{制約なしエントロピー双対}
注意~\ref{rem-uncons-dual}と同様に、$\int_\X \d\mu=\int_\Y\d\nu = 1$の場合、代替的な双対定式化を考えることができる
\eql{\label{eq-dual-entropic-alt}
	\usup{(\f,\g) \in \Cc(\X)\times\Cc(\Y)} \int_\X \f \d\al + \int_\Y \g \d\be
		+ \smine( c - f \oplus g ),
}
これは~\eqref{eq-dual-entropic}と同じ最適値を達成する。
%
~\eqref{eq-defn-softmin}と同様に、ソフト最小値（ここでは$\X \times \Y$上）は次のように定義される
\eq{
	\foralls S \in \Cc(\X\times\Y), \quad
	\smine S \eqdef
	- \varepsilon \int_{\X\times\Y} e^{ \frac{-S(x,y)}{\varepsilon} } \d\al(x)\d\be(y)
}
（これは$(\al,\be)$に依存することに注意する）。
%
$\epsilon \rightarrow 0$のとき、$\smine \rightarrow \min$であり、非正則化かつ制約なし定式化~\eqref{eq-dual-uncons}で使われるものに帰着する。
%
~\eqref{eq-dual-entropic}と~\eqref{eq-dual-entropic-alt}はともに制約なし問題であるが、
~\eqref{eq-dual-entropic-alt}の主な利点は条件が良いことである。
汎関数のHesse行列が$\epsilon$で一様に有界であるという意味においてである。
このような条件数の改善を得るもう1つの方法は半双対問題を考えることである；
\S\ref{sec-semi-discr-entropy}、特に注意~\ref{rem-second-order-smoothness}を参照。
%
この代替的な双対定式化の欠点は、対数の存在により\S\ref{sec-sgd}で詳述される確率的最適化手法の使用が妨げられることである；特に注意~\ref{rem-sgd-cont-cont}を参照。
\end{rem2}
%%%%%%%

%%%%%%%%%%%%%%%%%%%%%%%%%%%%%%%%%%%%%%%%%%%%%%%%%%%%%%%%%%%%%%%%%%%%%%%%%%
\section{最適輸送コストの正則化近似}
\label{sec-regularized-cost}

エントロピー双対~\eqref{eq-dual-formulation}は滑らかな制約なし凹最大化問題であり、以下の命題で詳述するように、もとのカントロビッチ双対~\eqref{eq-dual}を近似する。\todoK{explain soft indicator constraint, transition to result below.}

\begin{prop}\label{prop-feasibility-dual}
~\eqref{eq-dual-formulation}の最適解の任意のペア$(\fD^\star,\gD^\star)$は、$(\fD^\star,\gD^\star) \in \PotentialsD(\C)$（~\eqref{eq-feasible-potential}で定義された実行可能カントロビッチポテンシャルの集合）を満たす。結果として、任意の$\varepsilon$に対して
\eq{
	\dotp{\fD^\star}{\a} + \dotp{\gD^\star}{\b} \leq  \MKD_\C(\a,\b).
}
\end{prop}
\begin{proof}
	~\eqref{eq-prim-dual-ent}における主-双対最適性条件と、$\P$が確率分布であり従ってすべての$i,j$に対して$\P_{i,j}\leq 1$であるという制約から、$\exp(-(\fD^\star_i+\gD^\star_j-\C_{i,j})/\varepsilon)\leq 1$が得られ、したがって$\fD^\star_i+\gD^\star_j\leq\C_{i,j}$となる。
\end{proof}

~\eqref{eq-regularized-discr}で定義される正則化輸送コスト$\MKD_\C^\varepsilon$の主な利点は、それが滑らかで凸であることであり、変分問題における損失関数としての統合に最適である（第~\ref{c-variational}章を参照）。

\begin{prop}\label{prop-convexity-dual}
$\MKD_\C^\varepsilon(\a,\b)$は$\epsilon\geq 0$に対して$\a$と$\b$の同時凸関数である。$\epsilon>0$のとき、その勾配は次に等しい
\eq{
	\nabla\MKD_\C^\varepsilon(\a,\b)= \begin{bmatrix} \fD^\star \\ \gD^\star \end{bmatrix},
}
ここで$\fD^\star$と$\gD^\star$は、座標の和が0になるように選ばれた式~\eqref{eq-dual-formulation}の最適解である。
\end{prop}

\todoK{
\begin{proof}
	\todo{TODO}
\end{proof}
}


\citep{CuturiSinkhorn}において、2つのヒストグラム間のWasserstein距離を近似するための下界と上界が提案された。これらの界はSinkhornアルゴリズムによって得られる解において主問題と双対問題の目的関数を評価することから成る。

\begin{defn}[Sinkhorn・ダイバージェンス]
	$\fD^\star$と$\gD^\star$を~\eqref{eq-dual-formulation}の最適解、$\P^\star$を~\eqref{eq-regularized-discr}の解とする。Wasserstein距離は以下の主Sinkhorn・ダイバージェンスと双対Sinkhorn・ダイバージェンスを用いて近似される：
\begin{align*}
	\SINKHORNP_\C^\varepsilon(\a,\b)& \eqdef \dotp{\C}{\P^\star} =  \dotp{e^{\frac{\fD^\star}{\varepsilon}}}{(\K\odot \C)e^{\frac{\gD^\star}{\varepsilon}}}, \\
	\SINKHORND_\C^\varepsilon(\a,\b) & \eqdef \dotp{\fD^\star}{\a} + \dotp{\gD^\star}{\b},
\end{align*}
ここで$\odot$は行列の成分ごとの積を表す。
\end{defn}

\begin{prop}\label{prop-sinkhorn-div}
次の関係が成り立つ：
\eq{
	\SINKHORND_\C^\varepsilon(\a,\b)\leq \MKD_\C^\varepsilon(\a,\b) \leq \SINKHORNP_\C^\varepsilon(\a,\b).
}
さらに
\eql{\label{eq-sinkh-div-gap}
	\SINKHORNP_\C^\varepsilon(\a,\b) - \SINKHORND_\C^\varepsilon(\a,\b) = \varepsilon ( \HD(\P^\star)+1 ) .
}
\end{prop}
\begin{proof}
式~\eqref{eq-sinkh-div-gap}は、主問題と双対問題が最適値において同じ値を持つこと（~\eqref{eq-dual-formulation}を参照）から、
\eq{
	\MKD_\C^\varepsilon(\a,\b) = \SINKHORNP_\C^\varepsilon(\a,\b) - \varepsilon \HD(\P^\star) =
	\SINKHORND_\C^\varepsilon(\a,\b) - \varepsilon \dotp{e^{\fD^\star/\varepsilon} }{ \K e^{\gD^\star/\varepsilon}}
}
と書くことで得られる。最終結果は$\dotp{e^{\fD^\star/\varepsilon} }{ \K e^{\gD^\star/\varepsilon}}=1$であることに注意することで得られる。後者は$\P^\star$のすべての成分の和を計算することに相当するからである。
\end{proof}


上で与えた関係式はOT距離を実際に界で挟む実用的な方法を示唆するが、実際にはSinkhornアルゴリズムの収束時にのみ有効であり、したがって実際には真に有用ではない。実際には、Sinkhorn反復は常にある精度の閾値に達した後で終了される。所定の反復回数$L$が設定され、最適解$\fD^\star$と$\gD^\star$の代わりに反復$\itL{\fD}$と$\itL{\gD}$を用いて$\SINKHORND_\C^\varepsilon$を評価すると、下界が回復される：式~\eqref{eq-lse-sinkh-1}と~\eqref{eq-lse-sinkh-2}に現れる表記法を用いて、$\MKD_\C^\varepsilon$の次の有限ステップ近似を導入する：
\eql{\label{eq-algorithmic-loss}
	\itL{\SINKHORND_\C}(\a,\b) \eqdef \dotp{\itL{\fD}}{\a} + \dotp{\itL{\gD}}{\b}.
}
この「アルゴリズム的」Sinkhorn汎関数は$L\geq 1$の場合に正則化コスト関数の下界を与える。


\begin{prop}[有限Sinkhorn・ダイバージェンス]\label{defn-sinkhorn-div}
次の関係が成り立つ：
\eq{
	\itL{\SINKHORND_\C}(\a,\b) \leq \MKD_\C^\varepsilon(\a,\b).
}
\end{prop}
\begin{proof}命題~\ref{prop-feasibility-dual}の証明と同様に、たった1回のSinkhorn反復の後でも、~\eqref{eq-slse-sinkh-1}と~\eqref{eq-slse-sinkh-2}に従って、$\itL{\fD}$と$\itL{\gD}$は成分$\exp(-(\itL{\fD}_i+\itL{\gD}_j-\C_{i,j})/\varepsilon)$を持つ行列の列和が$\b$であり、したがってその各成分は1以下であるという事実を利用する。これにより$(\itL{\fD}_i,\itL{\gD})$の双対実行可能性が得られる。
\end{proof}

\begin{rem}[Sinkhorn反復の主問題における非実行可能性]
\eqref{eq-sink-matrix}で提供される主反復は主問題において実行可能ではないことに注意する。なぜなら定義上、これらの反復は収束時に周辺制約を満たすよう設計されているからである。したがって$\P^{(2L+1)}$は実行可能ではないため、$\dotp{\C}{\P^{(2L+1)}}$を$\MKD_\C(\a,\b)$の近似と見なすことは妥当ではない。注意~\ref{rem-complexity-rounding}で述べた~\citet{altschuler2017near}の丸めスキームを用いることで、$\MKD_\C^\varepsilon(\a,\b)$の上界を得ることができる。さらに、Sinkhornアルゴリズムと同様に~\citep{NIPS2018_8184}、複数のヒストグラムのペアに対して行列演算を用いて並列に便利に計算できる。
\end{rem}


\begin{rem}[有限双対Sinkhorn・ダイバージェンスの非凸性] ~\eqref{eq-dual-formulation}の正則化された式$\MKD_\C^\varepsilon$とは異なり、有限Sinkhorn・ダイバージェンス$\itL{\SINKHORND_\C}(\a,\b)$は一般にその引数の凸関数\emph{ではない}（数値的に容易に確認できる）。しかし$\itL{\SINKHORND_\C}(\a,\b)$は微分可能な関数であり、自動微分技法を用いて（注意~\ref{rem-auto-diff}を参照）任意の引数、特に$\C,\a$または$\b$に関して微分できる。
\end{rem}
\todoK{ automatically differentiated using the following recursion:\todo{write recursion}.}




%%%%%%%%%%%%%%%%%%%%%%%%%%%%%%%%%%%%%%%%%%%
\section{一般化Sinkhorn}
\label{sec-generalized}

正則化OT問題~\eqref{eq-regularized-discr}は、次の形の構造化凸最適化問題の特殊ケースである
\eql{\label{eq-generalized-ot-regul}
	\umin{\P}
		\sum_{i,j} \C_{i,j} \P_{i,j} - \varepsilon \HD(\P) + F( \P\ones_m ) + G( \transp{\P}\ones_n ).
}
実際、$F=\iota_{\{\a\}}$、$G=\iota_{\{\b\}}$と定義すると、閉凸集合$\Cc$の指示関数は
\eql{\label{eq-iota-function}
	\iota_{\Cc}(x) = \choice{
		0 \qifq x \in \Cc, \\
		+\infty \quad\text{otherwise},
	}
}
であり、$\Couplings(\a,\b)$を定義するハード周辺制約が復元される。
%
命題~\ref{prop-regularized-primal}の証明はこのより一般的な問題~\eqref{eq-generalized-ot-regul}にも適用されるため、~\eqref{eq-generalized-ot-regul}の一意の解も~\eqref{eq-scaling-form}の形を持つ。

~\citep{2015-Peyre-siims,FrognerNIPS,2016-chizat-sinkhorn,karlsson2016generalized}で示されたように、Sinkhorn反復~\eqref{eq-sinkhorn}はこの問題に拡張でき、次のように書ける
\eql{\label{eq-gen-sinkh}
	\uD \leftarrow \frac{\Prox_{F}^{\KLD}(\K \vD)}{\K \vD}
	\qandq
	\vD \leftarrow \frac{\Prox_{G}^{\KLD}(\transp{\K}\uD)}{\transp{\K}\uD},
}
ここで$\KLD$ダイバージェンスに対する近接作用素は
\eql{\label{eq-prox-kl}
	\foralls \uD \in \RR_+^N, \quad
	\Prox_{F}^{\KLD}(\uD) = \uargmin{\uD' \in \RR_+^N}	\KLD(\uD'|\uD) + F(\uD').
}
一部の関数$F,G$に対しては反復~\eqref{eq-gen-sinkh}の線形収束率を証明することが可能であり、これらのスキームは任意の測度に一般化できる；詳細は~\citep{2016-chizat-sinkhorn}を参照。

反復~\eqref{eq-gen-sinkh}は、$\Prox_{F}^{\KLD}$と$\Prox_{G}^{\KLD}$が閉じた形で、または非常に効率的に計算できる場合に興味深い。これは特に$F(\uD) = \sum_i F_i(\uD_i)$の形の分離可能関数の場合であり、このとき
\eq{
	\Prox_{F}^{\KLD}(\uD) = \pa{ \Prox_{F_i}^{\KLD}(\uD_i) }_i.
}
各$\Prox_{F_i}^{\KLD}$の計算は、スカラー最適化問題であるため通常単純である。
%
初期のSinkhornアルゴリズムと同様に、対数領域での計算を用いて計算を安定化することも可能であることに注意する~\citep{2016-chizat-sinkhorn}。

このアルゴリズムはOTの様々な一般化の解を近似するために使用でき、特に~\eqref{eq-unbalanced-pbm}の形の非均衡OT問題（\S\ref{sec-unbalanced}、特に反復~\eqref{eq-iterate-gen-sinkh}を参照）および~\eqref{eq-grad-flow-discr}の形の勾配流問題（\S\ref{sec-grad-flows}を参照）に使用できる。


%%%
\begin{rem1}{双対性とLegendre変換}
~\eqref{eq-generalized-ot-regul}の双対問題は次のようになる
\eql{\label{eq-dual-generalized}
	\umax{\fD,\gD}
		- F^*( \fD ) -G^*(\gD)
		- \varepsilon\sum_{i,j} e^{ \frac{\fD_i+\gD_j-\C_{i,j}}{\varepsilon} }
}
したがって$(\uD,\vD) = (e^{\fD/\varepsilon},e^{\gD/\varepsilon})$が~\eqref{eq-scaling-form}に現れる対応するスケーリングである。
ここで$F^*$と$G^*$はFenchel-Legendre共役であり、次のように定義される凸関数である
\eql{\label{eq-legendre}
		\foralls \fD \in \RR^n, \quad
		F^*(\fD) \eqdef \umax{\a \in \RR^n} \dotp{\fD}{\a} - F(\a).
}
%
一般化Sinkhorn反復~\eqref{eq-gen-sinkh}はDykstraアルゴリズム~\citep{Dykstra83,Dykstra85}（Bregman・ダイバージェンスに拡張されたもの~\citep{bauschke-lewis,CensorReich-Dykstra}；注意~\ref{rem-bregman}も参照）の特殊ケースであり、双対問題~\eqref{eq-dual-generalized}に対する交互最大化スキームである。
\end{rem1}
%%%

定式化~\eqref{eq-generalized-ot-regul}は2つ以上の関数および単一のカップリングを超えてさらに一般化できる；詳細は~\citep{2016-chizat-sinkhorn}を参照。これには特殊ケースとして、\S\ref{sec-multimarginal}で詳述する多周辺問題に対するSinkhornアルゴリズム~\eqref{eq-sinkh-multimarg}が含まれる。
%
正則化重心問題~\eqref{eq-entropic-bary}をこのように書き換えることも可能であり、反復~\eqref{eq-sinkhorn-bary}は実際にこの一般化Sinkhornの特殊ケースである。


\if 0 %% REMOVED %%%
\eq{
	\umin{(\P_s)_{s=1}^S}
		\sum_{i,j,s} \C_{s,i,j} \P_{s,i,j} - \varepsilon \HD(\P_s) +
		\sum_{k} F_k( (\P_s\ones_m)_s, (\transp{\P_s}\ones_n)_s )
}
where $\{\C_s\}_s$ are given cost matrices.

The solution has the form
\eql{\label{eq-scaling-form-multi}
	\foralls (s,i,j) \in \range{1,S} \times \range{n} \times \range{m}, \quad
	\P_{s,i,j} =  \K_{s,i,j} \prod_{k} \uD_{k,i} \prod_\ell \vD_{k,j}
}
where $\K_s=e^{-\C_s/\varepsilon}$ is a Gibbs kernel, and $(\uD_k,\vD_\ell)_{k,\ell}$ are scaling vectors that need to be computed using Generalized Sinkhorn iterations similar to~\eqref{eq-gen-sinkh}~\citep{2016-chizat-sinkhorn}.
%
In the case where the functions $F_s$ are indicator functions, one retrieves the Sinkhorn algorithm~\eqref{eq-sinkh-multimarg} for the multi-marginal problem, as detailed in~\S\ref{sec-multimarginal}.
%
It is also possible to rewrite the regularized barycenter problem~\eqref{eq-entropic-bary} this way, and the iterations~\eqref{eq-sinkhorn-bary} are in fact a special case of this generalized Sinkhorn.

\fi %%%%

\todoK{Write down the generalized iterations}

%These iterations read
%\eq{
%	\foralls k, \quad
%	\uD_s \leftarrow \frac{\Prox_{F_k}^{\KLD}( \tilde \uD_k \K \vD }{ \tilde \uD_k }
%	\qwhereq
%	\prod \uD_ \K \vD
%}
%\eq{
%	\vD \leftarrow \frac{\Prox_{G}^{\KLD}(\transp{\K}\uD)}{\transp{\K}\uD}.
%}
%where we defined the Gibbs kernel associated to the cost matrix $\C$ as
%\eq{
%	\K_{i,j} \eqdef e^{-\frac{\C_{i,j}}{\varepsilon}}
%}
